%%%%%%%%%%%%%%%%%%%%%%%%%%%%%%%%%%%%%%%%%%%%%%%%%%%%%%%%%%%%%%%%%%%%%%%%%%%%%
%                                                                           %
%   tuProlog Documentation                                                  %
%                                                                           %
%   A.Ricci, 5/12/00                                                        %
%   E.Denti, 31/5/01                                                        %
%   A.Ricci, 21/6/01                                                        %
%   A.Ricci, 13/9/01                                                        %
%   A.Ricci, 05/7/02                                                        %
%                                                                           %
%%%%%%%%%%%%%%%%%%%%%%%%%%%%%%%%%%%%%%%%%%%%%%%%%%%%%%%%%%%%%%%%%%%%%%%%%%%%%

\documentclass[11pt]{report}
%%%%%%%%%%%%%%%%%%%%%%%%%%%%%%%%%%%%%%%%%%%%%%%%%%%%%%%%%%%%%%%%%%%%%%%%%
 \newif\ifpdf
 \ifx\pdfoutput\undefined
 \pdffalse % we are not running PDFLaTeX
 \else
 \pdfoutput=1 % we are running PDFLaTeX
 \pdftrue
 \fi
%%%%%%%%%%%%%%%
 \ifpdf
 \usepackage[pdftex]{graphicx}
 \else
 \usepackage{graphicx}
 \fi
%%%%%%%%%%%%%%%
 \ifpdf
 \DeclareGraphicsExtensions{.eps, .pdf, .jpg, .tif}
 \else
 \DeclareGraphicsExtensions{.eps, .jpg}
 \fi
%%%%%%%%%%%%%%%
%\usepackage{harvard}
%%%%%%%%%%%%%%%
%\DeclareGraphicsRule{.gif}{bmp}{}{}% declare GIF filename extension
%\DeclareGraphicsRule{.jpg}{bmp}{}{}% declare WMF filename extension
%\DeclareGraphicsRule{.wmf}{bmp}{}{}% declare WMF filename extension
%%%%%%%%%%%%%%%%%%%%%%%%%%%%%%%%%%%%%%%%%%%%%%%%%%%%%%%%%%%%%%%%%%%%%%%%%
% -- reducing margins
 %\setlength{\oddsidemargin}{0pt}
 %\setlength{\evensidemargin}{0pt}
 %\setlength{\textwidth}{16cm}
 %
 %\setlength{\topmargin}{-0.4in}
 %\setlength{\textheight}{9.6in}
%%%%%%%%%%%%%%%%%%%%%%%%%%%%%%%%%%%%%%%%%%%%%%%%%%%%%%%%%%%%%%%%%%%%%%%%%
\newcommand\xa[1]{\appendixname~\ref{app:#1}}
\newcommand\labelsec[1]{\label{sec:#1}}
\newcommand\xs[1]{\sectionname~\ref{sec:#1}}
\newcommand\xsp[1]{\sectionname~\ref{sec:#1} \onpagename~\pageref{sec:#1}}
\newcommand\labelssec[1]{\label{ssec:#1}}
\newcommand\xss[1]{\subsectionname~\ref{ssec:#1}}
\newcommand\xssp[1]{\subsectionname~\ref{ssec:#1} \onpagename~\pageref{ssec:#1}}
\newcommand\labelsssec[1]{\label{sssec:#1}}
\newcommand\xsss[1]{\subsectionname~\ref{sssec:#1}}
\newcommand\xsssp[1]{\subsectionname~\ref{sssec:#1} \onpagename~\pageref{sssec:#1}}
\newcommand\labelfig[1]{\label{fig:#1}}
\newcommand\xf[1]{\figurename~\ref{fig:#1}}
\newcommand\xff[2]{\figurenames~\ref{fig:#1}~and~\ref{fig:#2}}
\newcommand\xfp[1]{\figurename~\ref{fig:#1} \onpagename~\pageref{fig:#1}}
\newcommand\labeltab[1]{\label{tb:#1}}
\newcommand\xt[1]{\tablename~\ref{tb:#1}}
\newcommand\xtt[2]{\tablenames~\ref{tb:#1}~and~\ref{ab:#2}}
\newcommand\xtp[1]{\tablename~\ref{tb:#1} \onpagename~\pageref{tb:#1}}
\newcommand\labelenum[1]{\label{enum:#1}}
\newcommand\xen[1]{(\ref{enum:#1})}
\newcommand\xenp[1]{(\ref{enum:#1}) \onpagename~\pageref{enum:#1}}
%******************************************************************************%
\newcommand\bti[1]{\texttt{\textbf{#1}}}
\newcommand\bt[1]{\texttt{#1}}
\newcommand\template[1]{\textit{Template: }\texttt{#1}}
\newcommand\ttit[1]{\texttt{\textit{#1}}}
%******************************************************************************%
\newcommand{\LIA}{\mbox{\textsf{LIA}}}
\newcommand{\aclt}{\mbox{$\mathcal{ACLT}$}}
\newcommand{\respect}{\mbox{\sf{{R}e{S}pec{T}}}}
\newcommand{\luce}{\mbox{\sf{{L}u{C}e}}}
\newcommand{\tucson}{\mbox{\sf{{T}u{CS}o{N}}}}
\newcommand{\alice}{\mbox{\sf{{aliCE}}}}
\newcommand{\tuprolog}{\mbox{{\sf{tu}}Prolog}}
%******************************************************************************%


\newcommand\version[1]{\mbox{Document revision: #1}}
\newcommand\approvedby[1]{\mbox{Approved by: #1}}
\newcommand\receivedby[1]{\mbox{Received by: #1}}
\newcommand\creationdate[1]{\mbox{Creation date: #1}}
\newcommand\lastchangesdate[1]{\mbox{Last Changes date: #1}}
\newcommand\noa[2]{\noindent\emph{Note of the author (#1): }#2\\\\}
\newcommand\logo{
    \begin{figure}[tp]
        \begin{center}
            \includegraphics[width=2cm]{images/logo}
        \end{center}
\end{figure}
}
%******************************************************************************
%
\title{{\huge{\bf{\tuprolog{} Developer's Guide\\\mbox{ }\\}}}
        \tuprolog{} version: 1.2.1\\\mbox{ }\\
{\small{
    \version{001}\\
    \creationdate{2002-09-17}\\
    \lastchangesdate{2004-03-19}\\
    % \receivedby{aricci}\\
    % \approvedby{aricci}\\
    }}
}

\author{ \mbox{ }\\DEIS, Universit\`{a} di Bologna a Cesena, Italy
}

\date{}

\begin{document}

\logo

\maketitle

\tableofcontents

%=======================================================================
\chapter{What is \tuprolog{}}
\label{what-is}
%=======================================================================

\tuprolog{} is an open-source, light-weight Prolog framework for distributed applications and infrastructures, released under the LGPL license, available from \url{http://tuprolog.apice.unibo.it}.

Originally developed in/upon Java, which still remains the main reference platform,
\tuprolog{} is currently available for several platforms/environments:
\begin{itemize}
  \item plain JavaSE;
  \item Eclipse plugin;
  \item Android;
  \item Microsoft .NET.
\end{itemize}
%
While they all share the same core and libraries, the latter features an \emph{ad hoc} library which extends the multi-paradigm approach to virtually any language available on the .NET platform (more on this in Section \xch{dotnet}).

Unlike most Prolog programming environments, aimed at providing a very efficient (yet monolithic) stand-alone Prolog system, \tuprolog{} is explicitly designed to be \emph{minimal}, dynamically \emph{configurable}, straightforwardly \emph{integrated} with Java and .NET so as to naturally support multi-paradigm/multi-language programming (MPP), and \emph{easily deployable}.

\textit{Minimality} means that its core contains only the Prolog engine essentials -- roughly speaking, the resolution engine and some related basic mechanisms -- for as little as 155KB: any other feature is implemented in \textit{libraries}.
%
So, each user can customize his/her prolog system to fit his/her own needs, and no more:
this is what we mean by \tuprolog{} \textit{configurability}---the necessary counterpart of minimality.

Libraries provide packages of predicates, functors and operators, and can be loaded and unloaded in a \tuprolog{} engine both statically and dynamically.
%
Several standard libraries are included in the \tuprolog{} distribution, and are loaded by default in the standard \tuprolog{} configuration; however, users can easily develop their own libraries either in several ways -- just pure Prolog, just pure Java\footnote{The .NET version of \tuprolog{} supports other languages available on the .NET platform: more on this topic in Section \xch{dotnet}}, or a mix of the two --, as we will discuss in \xch{howto-develop-libraries}.

\textit{Multi-paradigm programming} is another key feature of \tuprolog{}.
%
In fact, the \tuprolog{} design was intentionally calibrated from the early stages to support a straightforward, pervasive, multi-language/multi-paradigm integration, so as to enable users to:
\begin{itemize}
  \item using any Java\footnote{For the .NET version: any .NET class, library, object, etc.} class, library, object \emph{directly from the Prolog code}
  (\xch{java-library}) with no need of pre-declarations, awkward syntax, etc., with full support of parameter passing from the two worlds, yet leaving the two languages and computational models totally separate so as to preserve \emph{a priori} their own semantics---thus bringing the power of the object-oriented platform (e.g. Java Swing, JDBC, etc) to the Prolog world for free;

  \item using any Prolog engine \emph{directly from the Java/.NET code} as one would
   do with any other Java libraries/.NET assemblies, again with full support of parameter passing from the two worlds in a non-intrusive, simple way that does not alter any semantics---thus bringing the power of logic programming into virtually \emph{any} Java/.NET application;

  \item augmenting Prolog by defining new libraries \xch{howto-develop-libraries})
  either in Prolog, or in the object-oriented language of the selected platform (again, with a straightforward, easy-to-use approach based on reflection which avoids any pre-declaration, language-to-language mapping, etc), or in a mix of both;

  \item augmenting Java\footnote{This feature is currently available only in the Java version: a suitable extension to the .NET platform is under study.} by defining new Java methods in Prolog (the so-called `P@J' framework---\xch{PJ}), which exploits reflection and type inference to provide the user with an easy-to-use way to implement Java methods declaratively.
\end{itemize}

Last but not least, \textit{easy deployability} means that the installation requirements are minimal, and that the installation procedure is in most cases\footnote{Exceptions are the Eclipse plugin and the Android versions, which need to be installed as required by the hosting platforms.} as simple as copying one archive to the desired folder.
%
Coherently, a Java-based installation requires only a suitable Java Virtual Machine, and `installing' is just copying a single JAR file somewhere---for as much as 474KB of disk usage (yes, minimality is not just a claim here).
%
Of course, other components can be added (documentation, extra libraries, sources..), but are not necessary for a standard everyday use.
%
The file size is quite similar for the Android platform -- the single APK archive is 234KB -- although an Android-compliant install is performed due to Android requirements.
%
The install process is also quite the same on the .NET platform, although the files are slightly larger.
%
The Eclipse platform also requires a different procedure, since plugin installation have to conform to the requirements of the Eclipse plugin manager: consequently, an update site was set up, where the \tuprolog{} plugin is available as an Eclipse feature. Due to these constraints, file size increases to 1.5MB.


In order to manage all these platforms in a uniform way, a suitable \emph{version numbering scheme} was recently introduced:
\begin{itemize}
  \item the fist two digits represent the engine version;
  \item the last (third) digit is platform-specific and accounts for version differences which do not impact on the Prolog engine -- that is, on the
      \tuprolog{} behaviour -- but simply on graphical aspects or platform-specific
      issues or bugs.
\end{itemize}
%
So, as long as the first two digits are the same, a \tuprolog{} application is guaranteed to behave identically on any supported platform.


Finally, \tuprolog{} also supports \textit{interoperability} with both Internet standard patterns (such as TCP/IP, RMI, CORBA) and coordination models and languages.
%
The latter aspect, in particular, is currently developed in the context of the \tucson{} coordination infrastructure \cite{tucson-aamas99,respect-scico2001}, which provides logic-based, programmable tuple spaces (called \emph{tuple centres}) as the coordination media for distributed processes and agents.\footnote{An alternative infrastructure, \luce{} \cite{luce-aamas2001}, developed the same approach in a location-unaware fashion: this infrastructure is currently no longer supported.}



%\include{developerGuideInstallation}
%=======================================================================
\chapter{Using \tuprolog{} from Java}
\label{java-api}
%=======================================================================

\section{Getting started}

Let's begin with your first Java program using \tuprolog{}.
%
{\small{
\begin{verbatim}
    import alice.tuprolog.*;

    public class Test2P {
        public static void main(String[] args) throws Exception {
            Prolog engine = new Prolog();
            SolveInfo info = engine.solve("append([1],[2,3],X).");
            System.out.println(info.getSolution());
        }
    }
\end{verbatim}
}}
\noindent In this first example a \tuprolog{} engine is
created, and asked to solve a query, provided in a textual form.
%
This query in the Java environment is equivalent to the query\\\\
%
{\small{\texttt{?- append([1],[2,3],X).\\\\}}}
%
\noindent in a classic Prolog environment, and accounts for
finding the list that is obtained by appending the list
\texttt{[2,3]} to the list \texttt{[1]} (\texttt{append} is
included in the theory provided by the
\texttt{alice.tuprolog.lib.BasicLibrary}, which is downloaded by
the default when instantiating the engine).
%

By properly compiling and executing this simple program,\footnote{Save the program in a file called \texttt{Test2P.java}, then compile it with
%
\texttt{javac -classpath tuprolog.jar Test2P.java}
%
and then execute it with
%
\texttt{java -cp .;tuprolog.jar Test2P.java}} the string
\texttt{append([1],[2,3],[1,2,3])} -- that is the solution of out
query -- will be displayed on the standard output.
%
%

\noindent Then, let's consider a little bit more involved example:

{\small{
\begin{verbatim}
public class Test2P {
    public static void main(String[] args) throws Exception {
        Prolog engine = new Prolog();
        SolveInfo info = engine.solve("append(X,Y,[1,2]).");
        while (info.isSuccess()) {
            System.out.println("solution: " + info.getSolution() +
                               " - bindings: " + info);
            if (engine.hasOpenAlternatives()) {
                info = engine.solveNext();
            } else {
                break;
            }
        }
    }
}
\end{verbatim}
}}

In this case, all the solutions of a query are retrieved and
displayed, with also the variable bindings:
{\small{\begin{verbatim}

solution: append([],[1,2],[1,2]) - bindings: Y / [1,2]  X / []
solution: append([1],[2],[1,2]) - bindings: Y / [2]  X / [1]
solution: append([1,2],[],[1,2]) - bindings: Y / []  X / [1,2]

\end{verbatim}
}}

\section{Basic Data Structures}

\noindent All Prolog data objects are mapped onto Java objects:
\texttt{Term} is the base class for Prolog untyped terms such as
atoms and compound terms (represented by the \texttt{Struct} class),
Prolog variables (\texttt{Var} class) and Prolog typed terms such
as numbers (\texttt{Int}, \texttt{Long}, \texttt{Float},
\texttt{Double} classes).
%
In particular:

%
\begin{itemize}
    %
    \item \texttt{Term} -- this abstract class represents a generic Prolog term, and it is the
    root base class of \tuprolog{} data structures, defining the basic common services, such
    as term comparison, term unification and so on.
    %
    It is worth noting that it is an abstract class, so no direct \texttt{Term}
    objects can be instantiated;
    %
    %
    %
    \item \texttt{Var} -- this class (derived from \texttt{Term})
    represents \tuprolog{} variables.
    %
    A variable can be anonymous, created by means of the default constructor, with no
    arguments, or identified by a name, that must starts with an upper case letter or an
    underscore;
    %
    \item \texttt{Struct} -- this class (derived from \texttt{Term})
    represents un-typed \tuprolog{} terms, such as atoms, lists and compound terms;
    %
    \texttt{Struct} objects are characterised by a functor name and a
    list of arguments (which are \texttt{Term}s themselves), while
    \texttt{Var} objects are labelled by a string representing the
    Prolog name.
    %
    Atoms are mapped as \texttt{Struct}s with functor with a name and
    no arguments;
    %
    Lists are mapped as \texttt{Struct} objects with functor
    \verb|'.'|, and two \texttt{Term} arguments (head and tail of the list);
    %
    lists can be also built directly by exploiting the 2-arguments constructor, with
    head and tail terms as arguments.
    %
    Empty list is constructed by means of the no-argument constructor of \texttt{Struct}
    (default constructor).
    %
    \item \texttt{Number} -- this abstract class represents typed, numerical Prolog terms,
        and it is the  base class of \tuprolog{} number classes;
    %
    \item \texttt{Int, Long, Double, Float} -- these classes (derived from \texttt{Number})
    represent  typed numerical \tuprolog{} terms, respectively integer, long, double and
    float numbers.
    %
    Following the Java conventions, the default type for integer number is \texttt{Int}
    (integer, not long, number), and for \texttt{Double} (and so double), for floating point
    number.
\end{itemize}


\noindent Some examples of term creation follow:
%
{\small{\begin{verbatim}

// constructs the atom vodka
Struct drink = new Struct("vodka");

// constructs the number 40
Term degree = new alice.tuprolog.Int(40);

// constructs the compound degree(vodka, 40)
Term drinkDegree = new Struct("degree",
                              new Struct("vodka"),
                              new Int(40));
// second way to constructs the compound degree(vodka,40)
Struct drinkDegree2 = new Struct("degree", drink, degree);

// constructs the compound temperature('Rome', 25.5)
Struct temperature = new Struct("temperature",
                                new Struct("Rome"),
                                new alice.tuprolog.Float(25.5));

// constructs the compound equals(X, X)
Struct t1 = new Struct("equals", new Var("X"), new Var("X"));
t1.resolveTerm();

// mother(John,Mary)
Struct father = new Struct(new Struct("John"), new Struct("Mary")));

// father(John, _)
Term  father = new Struct(new Struct("John"), new Var());

// p(1, _, q(Y, 3.03, 'Hotel'))
Term  t2 = new Struct("p",
                      new Int(1),
                      new Var(),
                      new Struct("q",
                                 new Var("Y"),
                                 new Float(3.03f),
                                 new Struct("Hotel")));

// The Long number 130373303303
Term t3 = new alice.tuprolog.Long(130373303303h);

// The double precision number 1.7625465346573
Term t4 = new alice.tuprolog.Double(1.7625465346573);

// an empty list
Struct empty = new Struct();

// the list [303]
Struct l = new Struct(new Int(303), new Struct());

// the list [1,2,apples]
Struct alist = new Struct(
                   new Int(1),
                   new Struct(
                       new Int(2),
                       new Struct("apples")));

// fruits([apple, orange | _ ])
Term list2 = new Struct("fruits", new Struct(
                                      new Struct("apple",
                                          new Struct("orange"),
                                          new Var())));

// complex_compound(1, _, q(Y, 3.03, 'Hotel', k(Y,X)), [303, 13, _, Y])
Term t5 = Term.parse(
     "complex_compound(1, _, q(Y, 3.03, 'Hotel', k(Y,X)), [303, 13, _, Y])"
);

\end{verbatim}}}
%
\noindent The name of the \tuprolog{} number classes
(\texttt{Int}, \texttt{Float}, \texttt{Long}, \texttt{Double})
follows the name of the primitive Java data type they represents.
%
Note that due to class name clashes (for instance between classes
\texttt{java.lang.Long} and \texttt{alice.tuprolog.Long}), it
could be necessary to use the full class name to identify
\tuprolog{} classes.
%

\section{Engine, Theories and Libraries}

\noindent Then, the other main classes that make \tuprolog{} Core
API concern \tuprolog{} engines, theories and libraries. In
particular:
%
\begin{itemize}
    \item \texttt{Prolog} -- this class represent \tuprolog{}
    engines.
    %
    This class provides a minimal interface that enables users
    to: \\
    %
    \indent{-- set/get the theory to be used for
    demonstrations;}\\
        %
    \indent{-- load/unload libraries;} \\
        %
    \indent{-- solve a goal, represented either by a \texttt{Term} object or by a
        textual representation (a \texttt{String} object) of a
        term.}\\
    %
    A \tuprolog{} engine can be instantiated either with some standard default
    libraries loaded, by means of the default constructor, or with
    a starting set of libraries, which can be empty, provided as
    argument to the constructor (see JavaDoc documentation for
    details).
    %
    Accordingly, a raw, very lightweight, \tuprolog{} engine can
    be created by specifying an empty set of library, providing
    natively a very small set of built-in primitives.


    \item \texttt{Theory} -- this class represent \tuprolog{}
    theories.
    %
    A theory is represented by a text, consisting of a series of
    clauses and/or directives, each followed by a dot and a
    whitespace character.
    %
    Instances of this class are built either from a textual representation,
    directly provided as a string or taken by any possible input
    stream, or from a list of terms representing Prolog clauses.
    %
    %
    %
    %
    %
    \item \texttt{Library} -- this class represents \tuprolog{}
    libraries;
    %
    A \texttt{tuprolog} engine can be dynamically extended by loading
    (and unloading) any number of libraries; each library can provide
    a specific set of of built-ins predicates, functors and a related
    theory.
    %
    A library can be loaded by means of the built-in by means of the method
    \texttt{loadLibrary} of the \tuprolog{} engine.
    %
    Some standard libraries are provided in the
    \texttt{alice.tuprolog.lib} package and loaded by the default
    instantiation of a \tuprolog{} engine:
    %
    \texttt{alice.tuprolog.lib.BasicLibrary}, providing basic and
    well-known Prolog built-ins, \texttt{alice.tuprolog.lib.IOLibrary}
    providing \textit{de facto} standard Prolog I/O predicates, \texttt{alice.tuprolog.lib.ISOLibrary}
    providing some ISO predicates/functors not directly provided
    by \texttt{BasicLibrary} and \texttt{IOLibrary}, and
    \texttt{alice.tuprolog.lib.JavaLibrary}, which enables the
    creation and usage of Java objects from \tuprolog{} programs,
    in particular enabling the reuse of any existing Java resources.
    %
    %
    \item \texttt{SolveInfo} -- this class represents the result of a
    demonstration and instances of these class are returned by the
    \texttt{solve} methods the \texttt{Prolog} engines;
    %
    in particular \texttt{SolveInfo} objects provide services to test the
    success of the demonstration (\texttt{isSuccess} method),
    to access to the term solution of the query
    (\texttt{getSolution} method)  and to access the list of the
    variable with their bindings.
\end{itemize}
%

Some notes about \tuprolog{} terms and the services they provide:
\begin{itemize}
%
\item the static \texttt{parse} method provides a quick way to get a
term from its string representation.
%
\item \tuprolog{} terms provides directly methods for unification and matching: \\
%
{\tt{\small{public boolean unify(Term t)}}}\\
%
{\tt{\small{public boolean match(Term t)}}}\\
%
Terms that have been subject to unification outside a
demonstration context (that is invoking directly these methods,
and not passing through the solving service of an engine) should
not be used then in queries to engines.
%
\item some services are provided to compare terms, according to the
Prolog rules, and to check their type;
%
in particular the standard Java method \texttt{equals} has the
same semantics of the method \texttt{isEqual} which follows the
Prolog comparison semantics.
%
\item some services makes it possible to copy a term as it is
or to get a renamed copy of the term (\texttt{copy} and
\texttt{getRenamedCopy});
%
it is worth noting that the design of \tuprolog{} promotes a
stateless usage of terms; in particular, it is good practice not
to reuse the same terms in different demonstration contexts, as
part of different queries.
%
\item the method \texttt{getTerm} is useful in the case of
variables, providing the term linked possibly considering all the
linking chain in the case of variables referring other variables.
%
\item when a term is created by means of the proper constructor,
consider as example: \\\\
%
{\tt{\scriptsize{Struct myTerm = new Struct("p", new Var("X"), new
Int(1), new Var("X"))}}}\\\\
%
it is \emph{not resolved}, in the sense that possible variable
terms with the same name in the term do not refer each other;
%
so in the example the first and the third argument of the compound
\texttt{myTerm} point to different variable objects.
%
A term is resolved the first time it is involved in a matching or
unification context.
%
\end{itemize}

\noindent Some notes about \tuprolog{} engines, theories,
libraries and the services they provide:

\begin{itemize}
%
\item \tuprolog{} engines support natively some
\emph{directives}, that can be defined by means of the :-/1
predicate in theory specification.
%
Directives are used to specify properties of clauses and of the
engine (\emph{solve/1}, \emph{initialization/1},
\emph{set\_prolog\_flag/1}, \emph{load\_library/1},
\emph{consult/1}), format and syntax of read-terms (\emph{op/3},
\emph{char\_conversion/2}).
%
\item \tuprolog{} engines support natively the dynamic definition
and management of \emph{flags} (or property), used to describe
some aspects of libraries and their built-ins.
%
A flag is identified by a name (an alphanumeric atom), a list of
possible values, a default value and a boolean value specifying if
the flag value can be modified.
%
\item \tuprolog{} engines are thread-safe. The methods that could
create problems in being used in a multi-threaded context are now
synchronised.
%
\item \tuprolog{} engines have no (static) dependencies with each
other, multiple engines can be created independently as simple
objects on the same Java virtual machine,  each with its own
configuration (theory and loaded libraries).
%
Moreover, accordingly to the design of \tuprolog{} system in
general, engines are very lightweight, making suitable the use of
multiple engines in the same execution context.
%
\item \tuprolog{} engines can be serialised and stored as a persistent
object or sent through the network.
%
This is true also for engines with pre-loaded standard libraries:
%
in the case that other libraries are loaded, these must be
serializable in order to have the engine serializable.
%
\end{itemize}

%
%

\section{Some more examples of \tuprolog{} usage}

\noindent Creation of an engine (with default libraries
pre-loaded):

{\tt\scriptsize{
\begin{verbatim}
    import alice.tuprolog.*;

    ...
    Prolog engine = new Prolog();
\end{verbatim} }}


\noindent Creation of an engine specifying only the
\texttt{BasicLibrary} as pre-loaded library:

{\tt\scriptsize{
\begin{verbatim}
    import alice.tuprolog.*;

    ...
    Prolog engine = new Prolog(new String[]{"alice.tuprolog.lib.BasicLibrary"});
\end{verbatim} }}

\noindent Creation and loading of a theory from a string:

{\tt\scriptsize{\begin{verbatim}

    String theoryText = "my_append([],X,X).\n" +
                        "my_append([X|L1],L2,[X|L3]) :- my_append(L1,L2,L3).\n";

    Theory theory = new Theory(theoryText);
    try {
        engine.setTheory(theory);
    } catch(InvalidTheoryException e) {
    }
\end{verbatim} }}

\noindent Creation and loading of a theory from an input stream:

{\tt\scriptsize{\begin{verbatim}

    Theory theory = new Theory(new FileInputStream("test.pl");
    try {
        engine.setTheory(theory);
    } catch(InvalidTheoryException e) {
    }
\end{verbatim} }}

\noindent Goal demonstration (provided as a string):

{\tt\scriptsize{\begin{verbatim}

    // ?- append(X,Y,[1,2,3]).
    try {
        SolveInfo info = engine.solve("append(X,Y,[1,2,3]).");
        Term solution = info.getSolution();
    } catch(MalformedGoalException mge) {
        ...
    } catch(NoSolutionException nse) {
        ...
    }
\end{verbatim} }}

\noindent Goal demonstration (provided as a Term):

{\tt\scriptsize{\begin{verbatim}

    try {
        Term goal = new Struct("p", new Int(1), new Var("X"));
        try {
            // ?- p(1,X).
            SolveInfo info = engine.solve(goal);
            Term solution = info.getSolution();
 
        } catch (NoSolutionException nse) {
        }
    } catch (InvalidVarNameException ivne) {
    }
\end{verbatim} }}

\noindent Getting another solution:

{\tt\scriptsize{\begin{verbatim}
    try {
        SolveInfo info = engine.solve(goal);
        info = engine.solveNext();
    } catch(NoMoreSolutionException e)
\end{verbatim} }}

\noindent Loading a library:

{\tt\scriptsize{\begin{verbatim}
    try {
        engine.loadLibrary('alice.tuprologx.lib.TucsonLibrary');
    } catch(InvalidLibraryException e) {
    }
\end{verbatim} }}

\noindent Here, a complete example of interaction with a
\tuprolog{} engine is shown (refer to the JavaDoc documentation for
details about interfaces):

{\tt\scriptsize{\begin{verbatim}

import alice.tuprolog.*; import java.io.*;

public class Test2P {
    public static void main (String args[]) {
        Prolog engine = new Prolog();
        try {

            // solving a goal
            SolveInfo info = engine.solve(new Struct("append",
                                              new Var("X"),
                                              new Var("Y"),
                                              new Struct(new Term[]{new Struct("hotel"),
                                                                    new Int(303),
                                                                    new Var()})));


            // note we could use strings:
            // SolveInfo info = engine.solve("append(X, Y, [hotel, 303, _]).");

            // test for demonsration success
            if (info.isSuccess()) {

                // acquire solution and substitution
                Term sol = info.getSolution();
                System.out.println("Solution: " + sol);

                System.out.println("Bindings: " + info);

                // open choice points?
                if (engine.hasOpenAlternatives()) {

                    // ask for another solution
                    info = engine.solveNext();

                    if (info.isSuccess()) {
                        System.out.println("An other substitution: " + info);
                    }
                }
            }

            // other frequent interactions

            // setting a new theory in the engine
            String theory = "p(X,Y) :- q(X), r(Y).\n" +
                            "q(1).\n" +
                            "r(1).\n" +
                            "r(2).\n";
            engine.setTheory(new Theory(theory));

            SolveInfo info2 = engine.solve("p(1,X).");
            System.out.println(info2);

            // retrieving the theory from a file
            FileOutputStream os=new FileOutputStream("test.pl");
            os.write(theory.getBytes());
            os.close();
            engine.setTheory(new Theory(new FileInputStream("test.pl")));
            info2 = engine.solve("p(X,X).");
            System.out.println(info2.getSolution());

        } catch (Exception ex) {
            ex.printStackTrace();
        }
    }
}
\end{verbatim}}}

With the program execution, the following string are displayed on
the standard output:

{\tt\small{

\begin{verbatim}
Solution: append([],[hotel,303,_],[hotel,303,_])
Bindings: Y /[hotel,303,_] X / []
An other substitution: Y / [303,_]  X / [hotel]
X / 1
p(1,1)
\end{verbatim}}}
%=======================================================================
\chapter{Accessing Java from \tuprolog{}}
\label{java-library}
%=======================================================================
One of the main advantages of \tuprolog{} open architecture is
that any Java component can be directly accessed and used from
Prolog, in a simple and effective way, by means of the
\texttt{JavaLibrary} library: this delivers all the power of
existing Java components and packages to \tuprolog{} sources.
%
In this way, all Java packages involving interaction (such as Swing,
JDBC, the socket package, RMI) are immediately available to increase
the interaction abilities of \tuprolog:
%
{``one library for all libraries''} is the basic motto.
%
%
\section{Mapping data structures}

Complete bi-directional mapping is provided between Java primitive
types and \tuprolog{} data types.
%
By default, \tuprolog{} integers are mapped into Java \texttt{int}
or \texttt{long} as appropriate, while \texttt{byte} and
\texttt{short} types are mapped into \tuprolog{}'s \texttt{Int}
instances. Only Java \texttt{double} numbers are used to map
\tuprolog{} reals, but \texttt{float} values returned as result of
method invocations or field accesses are handled properly anyway,
without any loss of information.
%
Boolean Java values are mapped into specific \tuprolog{}
\texttt{Term} constants.
%
Java \texttt{char}s are mapped into Prolog atoms, but atoms are
mapped into Java \texttt{String}s by default.
%
The \emph{any} variable (\_) can be used to specify the Java
\texttt{null} value.

%---------------------------------------------------------------
\section{General predicates description}
%---------------------------------------------------------------

\begin{figure}
\caption{A sample Java class (a counter) used to explain JavaLibrary predicates behaviour.
\labelfig{jreflect-example}}
\begin{verbatim}
public class Counter {
    public String name;
    private long value = 0;

    public Counter() {}
    public Counter(String aName) { name = aName; }

    public void setValue(long val) { value=val; }
    public long getValue() { return value; }
    public void inc() { value++; }

    static public String getVersion() { return "1.0"; }
}
\end{verbatim}
\end{figure}

The library offers the following predicates:
%
\begin{enumerate}
  \renewcommand\labelenumi{\it(\roman{enumi})}
  %
  \item the \texttt{java\_object/3} predicate is used to create a new Java
        object of the specified class, according to the syntax:
        %
        \begin{center}
        \texttt{java\_object(\textit{ClassName},
                             \textit{ArgumentList},
                             \textit{ObjectRef})}
        \end{center}
        %
        \texttt{\textit{ClassName}} is a Prolog atom bound to the name of the
        proper Java class (e.g. \verb|'Counter'|, \verb|'java.io.FileInputStream'|),
        while the parameter \texttt{\textit{ArgumentList}} is a Prolog list used to supply
        the required arguments to the class
        constructor: the empty list matches the default constructor.
        %
        %%%%%% RICCI 020202
        Also Java arrays can be instantiated, by appending
        \texttt{[]} at the end of the \texttt{\textit{ClassName}}
        string.
        %%%%%%
        %
        The reference to the newly-created object is bound to \texttt{\textit{ObjectRef}},
        which is typically a ground Prolog term; alternatively, an unbound term
        may be used, in which case the term is bound to an automatically-generated
        %%%%%% RICCI 020202
        Prolog atom \verb|'$obj_N'|, where \texttt{N} is a progressive integer.
        %%%%%%
        %
        In both cases, these atoms are interpreted as object references --
        and therefore used to operate on the Java object from Prolog -- \textit{only}
        in the context of \texttt{JavaLibrary}'s predicates.
        %
        %%%%%% RICCI 020202
        %
        The predicate fails whenever \textit{ClassName} does not identify a valid Java class,
        or the constructor does not exists, or arguments in
        \texttt{\textit{ArgumentList}} are not ground, or \textit{ObjectRef}
        already identifies an object in the system.
        %
        %%%%%%

        According to the default behaviour of \texttt{java\_object},
        when a ground term is bound to a Java object by means of the predicate,
        the binding is kept for the full time of the demonstration
        (even in the case of backtracking).
        %
        This behaviour can be changed, getting the bindings
        created by the \texttt{java\_object} undone by
        backtracking, by changing the value of the flag \texttt{java\_object\_backtrackable}
        to \texttt{true} (the default is \texttt{false}).



  \item the \texttt{<-/2} predicate is used to invoke a method on a Java
        object according to a send-message pattern:
        %
        \begin{center}
        \texttt{\textit{ObjectRef} <- \textit{MethodName}(\textit{Arguments})}

        \texttt{\textit{ObjectRef} <- \textit{MethodName}(\textit{Arguments})
                returns \textit{Term}}
        \end{center}
        %
        \texttt{\textit{ObjectRef}} is an atom interpreted as a Java object
        reference as explained above, while \texttt{\textit{MethodName}}
        is the Java name of the method to be invoked, along with its
        \texttt{\textit{Arguments}}.
        %
        The \texttt{returns} keyword is used to retrieve the value returned
        from non-void Java methods and bind it to a Prolog term: if
        the type of the returned value can be mapped onto a primitive Prolog
        data type (a number or a string), \texttt{\textit{Term}} is unified
        with the corresponding Prolog value; if, instead, it is a Java object
        other than the ones above, \texttt{\textit{Term}} is handled
        as \texttt{\textit{ObjectRef}} in the case of \texttt{java\_object/3}.
        %
        %%%%%% RICCI 020202
        %
        % - static methods access
        %
        Static methods can be invoked using the compound
        term \texttt{class(\textit{ClassName})} in the place
        of \texttt{\textit{ObjectRef}}.
        %
        % - accennare al most specific method?
        %
        %%%%%%
        %%%%%% RICCI 020202
        %
        % - the predicates fails if...
        %
        If \textit{MethodName} does not identify a valid method for the object (class),
        or arguments in \texttt{\textit{ArgumentList}} are not
        valid (because of a wrong signature or not ground values) the predicate fails.
        %%%%%%

  \item the \texttt{.} infix operator is used, in conjunction with the \texttt{set}
        / \texttt{get} pseudo-method pair, to access the public fields of a Java
        object.
        %
        The syntax is based on the following constructs:
        %
        \begin{center}
        \tt
        \textit{ObjectRef} . \textit{Field} <- set(\textit{GroundTerm})\\
        \textit{ObjectRef} . \textit{Field} <- get(\textit{Term})\\
        \end{center}
        %
        As usual, \texttt{\textit{ObjectRef}} is the Prolog identifier for
        a Java object.
        %
        The first construct set the public field \texttt{\textit{Field}}
        to the specified \texttt{\textit{GroundTerm}}, which may be either
        a value of a primitive data type, or a reference to an existing
        object: if \texttt{\textit{GroundTerm}} is not ground, the infix
        predicate fails.
        %
        The second construct retrieves the value of the public field
        \texttt{\textit{Field}}, where \texttt{\textit{Term}} is handled
        once again as \texttt{\textit{ObjectRef}} in the case of
        \texttt{java\_object/3}.
        %
        %%%%%% RICCI 020202
        %
        % - static class access
        %
        As for methods, static fields of classes can be accessed using the compound
        term \texttt{class(\textit{ClassName})} in the place
        of \texttt{\textit{ObjectRef}}.
        %
        % - accesso ad array
        Some helper predicates are provided to access Java array
        elements:\\
        \texttt{java\_array\_set(\textit{ArrayRef}, \textit{Index}, \textit{Object})}\\
        \texttt{java\_array\_set\_\textit{\emph{Basic Type}}(\textit{ArrayRef}, \textit{Index}, \textit{Value})}\\
        to set elements,\\
        \texttt{java\_array\_get(\textit{ArrayRef}, \textit{Index}, \textit{Object})}\\
        \texttt{java\_array\_get\_\textit{\emph{Basic Type}}(\textit{ArrayRef}, \textit{Index}, \textit{Value})}\\
        to get elements,\\
        \texttt{java\_array\_length(\textit{ArrayObject}, \textit{Size})}
        to get the array length.\\
        %
        %%%%%%
        It is worth to point out that the \texttt{set} and \texttt{get} formal
        pseudo-methods above are \textit{not} methods of some class, but just
        part of the construct of the \texttt{.} infix operator, according to
        a JavaBeans-like approach.

  \item the \texttt{as} infix operator is used to explicitly specify (i.e., cast)
        method argument types:
        %
        \begin{center}
        \texttt{\textit{ObjectRef} as \textit{ClassName}}
        \end{center}
        %
        By writing so, the object represented by \texttt{\textit{ObjectRef}} is
        considered to belong to class \texttt{\textit{Classname}}: both
        \texttt{\textit{ObjectRef}} and \texttt{\textit{Classname}} have
        the usual meaning explained above.
        %
        %%%%%% RICCI 020202
        The operator works also with primitive Java types, specified
        as \texttt{\textit{Classname}} (for instance, \texttt{myNumber \textit{as int}}).
        %%%%%%
        %
        The purpose of this predicate is both to provide methods with the
        exact Java types required, and to solve possible overloading conflicts
        a-priori.
        %
        %%%%%% RICCI 020202
        %
        %In particular, \texttt{as} is needed when invoking a method whose a
        %formal argument is of type \texttt{\textit{T}} (e.g., \texttt{JButton})
        %with an actual object argument whose type is a subclass of
        %\texttt{\textit{T}} (say, \texttt{myButton}) -- that is, when upcasting
        %is needed to let Java identify the method signature unambiguously.
        %%%%%%

  %%%%%% RICCI 020202
  %
  \item The \texttt{java\_class/4} predicate makes it possible
        to create and load a new Java class from a source text provided as an
        argument, thus supporting \textit{dynamic compilation} of Java
        classes:
        %
        %
        \begin{center}
        \texttt{java\_class(\textit{SourceText},
                            \textit{FullClassName},
                            \textit{ClassPathList},
                            \textit{ObjectRef})}
        \end{center}
        %
        \texttt{\textit{SourceText}} is a string representing the
        text source of the Java class, \texttt{\textit{FullClassName}}
        is the full Java class name, and \texttt{\textit{ClassPathList}}
        is a (possibly empty) Prolog list of class paths that may
        be required for a successful dynamic compilation
        of this class.
        %
        \texttt{\textit{ObjectRef}} is a reference to an instance of the
        class \texttt{java.lang.Class} that represents the newly-created class.
        %
        The predicate fails whenever \texttt{\textit{SourceText}} contains errors,
        or the class cannot be located in the package hierarchy
        as specified, or \texttt{\textit{ObjectRef}} already identifies an object
        in the system.
  %
  %%%%%%

\end{enumerate}
%
%%%%%% RICCI 020202
%
\noindent Generally, exceptions thrown by method or constructor
calls cannot be explicitly managed and cause the failure of the
related predicate.
%

To taste the flavour of \texttt{JavaLibrary}, let us consider the
example below (refer to \xf{jreflect-example} for \texttt{Counter}
class definition):

%
{\small
\begin{verbatim}
    ?-  java_object('Counter', ['MyCounter'], myCounter),
        myCounter <- setValue(5),
        myCounter <- inc,
        myCounter <- getValue returns Value,
        write(X),

        class('Counter') <- getVersion return Version,

        myCounter.name <- get(Name),
        class('java.lang.System') . out <- get(Out),
        Out <- println(Name),

        myCounter.name <- set('MyCounter2'),

        java_object('Counter[]', [10], ArrayCounters),
        java_array_set(ArrayCounters, 0, myCounter).
\end{verbatim}}
%
\noindent Here, a \texttt{Counter} object is created providing the
\texttt{MyCounter} name as constructor argument: the reference to
the new object is bound to the Prolog atom \texttt{myCounter}.
%
This reference is then used for method invocation via the
\texttt{<-} operator, calling the \texttt{setValue(5)} method
(which is void and therefore returns nothing) first, incrementing
the counter (no arguments are specified) and invoking the
\texttt{getValue} method just after.
%
Since \texttt{getValue} returns an integer value, the
\texttt{returns} operator retrieves the method result (hopefully,
5) and binds it to the \texttt{X} Prolog variable, which is
printed via the Prolog \texttt{write/1} predicate.
%
Of course, if the Prolog variable \texttt{X} is already bound to
5, the predicate succeeds as well, while fails if \texttt{X} is
bound to anything else.
%
%%%%%%
Then, the static method \texttt{getVersion} is called, retrieving
the version of the class \texttt{Counter}, and printed using the
method \texttt{println} provided by the static \texttt{out} field
in the \texttt{java.lang.System} class.
%
The \texttt{name} public field of \texttt{myCounter} object is
then accessed, setting the \texttt{MyCounter2} value.
%
Finally, an array of 10 counters is created, and the
\texttt{myCounter} object assigned to its first element.
%%%%%%
%

The key point here is that the only requirement for this example
to run is the presence of the \texttt{Counter.class} file in the
proper position in the file system, according to Java naming
conventions: no other auxiliary information is needed -- no
headers, no pre-compilations, etc.
%
This enables the seamless reuse and exploitation of the large
amount of available Java libraries and resources, starting from
the standard ones, such as Swing to manage GUI components, JDBC to
access databases, RMI and CORBA for distributed computing, and so
on.
%
\xf{jexamples-swing} shows an example, where Java Swing API is
exploited to graphically choose a file from Prolog: a Swing
\texttt{JFileChooser} dialog is instantiated and bound to the
Prolog variable \texttt{Dialog} (a univocal Prolog atom of the form
\verb|'$obj_N'|, to be
used as the object reference, is automatically generated and bounded
to the variable) which
is then used to invoke methods \texttt{showOpenDialog} and
\texttt{getSelectedFile} of \texttt{JFileChooser}'s interface.
%
Further examples about exploiting standard Java libraries from
\tuprolog{}
%can be found in the Appendix A.
can be found in \cite{tuprolog-padl2001}.

\begin{figure}
\caption{Using a Swing componen from a \tuprolog{} program. Note the \texttt{\_} Prolog value used to represent the Java \texttt{null} value.
\labelfig{jexamples-swing}}
\begin{verbatim}
test_open_file_dialog(FileName) :-
    java_object('javax.swing.JFileChooser', [], Dialog),
    Dialog <- showOpenDialog(_),
    Dialog <- getSelectedFile returns File,
    File <- getName returns FileName.
\end{verbatim}
\end{figure}

Besides the Prolog predicates, \texttt{JavaLibrary} embeds the
\texttt{register} function, which, unlike the previous
functionalities, is to be used on the Java side.
%
Its purpose is to associate an existing Java object
\texttt{\textit{obj}} to a Prolog identifier
\texttt{\textit{ObjectRef}}, according to the syntax:
%
\begin{center}
 \small\tt
 boolean register(Struct \textit{ObjectRef}, Object \textit{obj})
    throws InvalidObjectIdException;\\
\end{center}
%
\texttt{\textit{ObjectRef}} is a ground term (otherwise an
exception is raised) that represents the Java object
\texttt{\textit{obj}} in the context of \texttt{JavaLibrary}'s
predicates: the function returns \texttt{false} if the object
represented by \texttt{\textit{obj}} is already registered under a
different \texttt{\textit{ObjectRef}}.
%
As an example of use, let us consider the following
case:\footnote{An
  explicit cast to \texttt{alice.tuprolog.lib.JavaLibrary} is needed because
  \texttt{loadLibrary} returns a reference to a generic
  \texttt{Library}, while the \texttt{register} primitive is defined in
  \texttt{JavaLibrary} only.}
%
{\small
\begin{verbatim}
Prolog core = new Prolog();
Library lib = core.loadLibrary("alice.tuprolog.lib.JavaLibrary");
((alice.tuprolog.lib.JavaLibrary)lib).register(new Struct("stdout"),
                                               System.out);
\end{verbatim}}
%
\noindent Here, the Java object \texttt{System.out} is registered
for use in \tuprolog{} under the name \texttt{stdout}.
%
So, within the scope of the \texttt{core} engine, a Prolog program
can now contain
\begin{verbatim}
stdout <- println('What a nice message!')
\end{verbatim}
as if \texttt{stdout} was a pre-defined \tuprolog{} identifier.

%---------------------------------------------------------------------
\section{Predicates}
%---------------------------------------------------------------------

\noindent Here follows a list of predicates implemented by this
library, grouped in categories corresponding to the functionalities
they provide.

%---------------------------------------------------------------------
\subsection{Method Invocation, Object and Class Creation}
%---------------------------------------------------------------------

\begin{itemize}
%
\item \bti{java\_object/3}\\
\noindent\bt{java\_object(ClassName, ArgList, ObjId)} is true iff
\bt{ClassName} is the full class name of a Java class available on
the local file system, \bt{ArgList} is a list of arguments that
can be meaningfully used to instantiate an object of the class,
and \bt{ObjId} can be used to reference such an object;
%
as a side effect, the Java object is created and the reference to
it is unified with \bt{ObjId}.
%
It is worth noting that \bt{ObjId} can be a Prolog variable (that
will be bound to a ground term) as well as a ground term (not a
number).\\
\template{java\_object(+full\_class\_name, +list, ?obj\_id)}
%
\item \bti{java\_object\_bt/3}\\
\noindent\bt{java\_object\_bt(ClassName, ArgList, ObjId)} has the same behaviour of \bt{java\_object/3}, but the binding that is established between the \bt{ObjId} term and the Java object is destroyed with backtracking.\\
\template{java\_object\_bt(+full\_class\_name, +list, ?obj\_id)}
%
\item \bti{destroy\_object/1}\\
\noindent\bt{destroy\_object(ObjId)} is true and as a side effect
the binding between \bt{ObjId} and a Java object,
possibly established, by previous predicates is destroyed.\\
\template{destroy\_object(@obj\_id)}
%
\item \bti{java\_class/4}\\
\noindent\bt{java\_class(ClassSourceText, FullClassName, ClassPathList, ObjId)}
is true iff \bt{ClassSouceText} is a source string describing a
valid Java class declaration, a class whose full name is
\bt{FullClassName}, according to the classes found in paths
listed in \bt{ClassPathList}, and \bt{ObjId} can be used as a
meaningful reference for a \texttt{java.lang.Class} object
representing that class;
%
as a side effect the described class is (possibly created and)
loaded and made available to the system.\\
\template{java\_class(@java\_source, @full\_class\_name, @list, ?obj\_id)}
%
\item \bti{java\_call/3}\\
\noindent\bt{java\_call(ObjId, MethodInfo, ObjIdResult)} is true iff
\bt{ObjId} is a ground term currently referencing a Java object,
which provides a method whose name is the functor name of the term
\bt{MethodInfo} and possible arguments the arguments of
\bt{MethodInfo} as a compound, and \bt{ObjIdResult} can be used as
a meaningful reference for the Java object that the method
possibly returns.
%
As a side effect the method is called on the Java object
referenced by the \bt{ObjId} and the object possibly returned by
the method invocation is referenced by the \bt{ObjIdResult} term.
%
The anonymous variable used as argument in the \bt{MethodInfo}
structure is interpreted as the Java \texttt{null} value.\\
\template{java\_call(@obj\_id, @method\_signature, ?obj\_id)}
%
\item \verb|'<-'/2|\\
\noindent\verb|'<-'(ObjId, MethodInfo)| is true iff \bt{ObjId} is
a ground term currently referencing a Java object, which provides a
method whose name is the functor name of the term \bt{MethodInfo}
and possible arguments the arguments of \bt{MethodInfo} as a
compound.
%
As a side effect the method is called on the Java object
referenced by the \bt{ObjId}.
%
The anonymous variable used as argument in the \bt{MethodInfo}
structure is interpreted as the Java \texttt{null} value.\\
\template{'<-'(@obj\_id, @method\_signature)}
%
\item \bti{return/2}\\
\noindent\verb|return('<-'(ObjId, MethodInfo), ObjIdResult)| is true
iff \bt{ObjId} is a ground term currently referencing a Java object,
which provides a method whose name is the functor name of the term
\bt{MethodInfo} and possible arguments the arguments of
\bt{MethodInfo} as a compound, and \bt{ObjIdResult} can be used as
a meaningful reference for the Java object that the method
possibly returns.
%
As a side effect the method is called on the Java object
referenced by the \bt{ObjId} and the object possibly returned by
the method invocation is referenced by the \bt{ObjIdResult} term.
%
The anonymous variable used as argument in the \bt{MethodInfo}
structure is interpreted as the Java \texttt{null} value.\\
%
It is worth noting that this predicate is equivalent to the
\texttt{java\_call} predicate.\\
\template{return('<-'(@obj\_id, @method\_signature), ?obj\_id)}
%
\end{itemize}

%---------------------------------------------------------------------
\subsection{Java Array Management}
%---------------------------------------------------------------------
\begin{itemize}
%
\item \bti{java\_array\_set/3}\\
\noindent\bt{java\_array\_set(ObjArrayId, Index, ObjId)} is true iff
\bt{ObjArrayId} is a ground term currently referencing a Java
array object, \bt{Index} is a valid index for the array and
\bt{ObjId} is a ground term currently referencing a Java object
that could inserted as an element of the array (according to Java
type rules).
%
As side effect, the object referenced by \bt{ObjId} is set in the
array referenced by \bt{ObjArrayId} in the position (starting from
0, following the Java convention) specified by \bt{Index}.
%
The anonymous variable used as \bt{ObjId} is interpreted as the
Java \texttt{null} value.
%
This predicate can be used for arrays of Java objects:
%
for arrays whose elements are Java primitive types (such as
\texttt{int}, \texttt{float}, etc.) the following predicates can
be used, with the same semantics of \bt{java\_array\_set} but
specifying directly the term to be set as a \tuprolog{} term
(according to the mapping described previously):\\
%
\mbox{~~~~}\bt{java\_array\_set\_int(ObjArrayId, Index, Integer)}\\
\mbox{~~~~}\bt{java\_array\_set\_short(ObjArrayId, Index, ShortInteger)}\\
\mbox{~~~~}\bt{java\_array\_set\_long(ObjArrayId, Index, LongInteger)}\\
\mbox{~~~~}\bt{java\_array\_set\_float(ObjArrayId, Index, Float)}\\
\mbox{~~~~}\bt{java\_array\_set\_double(ObjArrayId, Index, Double)}\\
\mbox{~~~~}\bt{java\_array\_set\_char(ObjArrayId, Index, Char)}\\
\mbox{~~~~}\bt{java\_array\_set\_byte(ObjArrayId, Index, Byte)}\\
\mbox{~~~~}\bt{java\_array\_set\_boolean(ObjArrayId, Index, Boolean)}\\
%
\template{java\_array\_set(@obj\_id, @positive\_integer, +obj\_id)}
%
%
%
\item \bti{java\_array\_get/3}\\
\noindent\bt{java\_array\_get(ObjArrayId, Index, ObjIdResult)} is
true iff \bt{ObjArrayId} is a ground term currently referencing a
Java array object, \bt{Index} is a valid index for the array, and
\bt{ObjIdResult} can be used as a meaningful reference for a Java
object contained in the array.
%
As a side effect, \bt{ObjIdResult} is unified with the reference to
the Java object of the array referenced by \bt{ObjArrayId} in the
\bt{Index} position.
%
This predicate can be used for arrays of Java objects:
%
for arrays whose elements are Java primitive types (such as
\texttt{int}, \texttt{float}, etc.) the following predicates can
be used, with the same semantics of \bt{java\_array\_get} but
binding directly the array element to a \tuprolog{} term
(according to the mapping described previously):\\
%
\mbox{~~~~}\bt{java\_array\_get\_int(ObjArrayId, Index, Integer)}\\
\mbox{~~~~}\bt{java\_array\_get\_short(ObjArrayId, Index, ShortInteger)}\\
\mbox{~~~~}\bt{java\_array\_get\_long(ObjArrayId, Index, LongInteger)}\\
\mbox{~~~~}\bt{java\_array\_get\_float(ObjArrayId, Index, Float)}\\
\mbox{~~~~}\bt{java\_array\_get\_double(ObjArrayId, Index, Double)}\\
\mbox{~~~~}\bt{java\_array\_get\_char(ObjArrayId, Index, Char)}\\
\mbox{~~~~}\bt{java\_array\_get\_byte(ObjArrayId, Index, Byte)}\\
\mbox{~~~~}\bt{java\_array\_get\_boolean(ObjArrayId, Index, Boolean)}\\
%
\template{java\_array\_get(@obj\_id, @positive\_integer, ?obj\_id)}
%
%
\item \bti{java\_array\_length/2}\\
\noindent\bt{java\_array\_length(ObjArrayId, ArrayLength)} is true
iff \bt{ArrayLength} is the length of the Java array referenced by
the term \bt{ObjArrayId}.\\
\template{java\_array\_length(@term, ?integer)}
%
\end{itemize}

%---------------------------------------------------------------------
\subsection{Helper Predicates}
%---------------------------------------------------------------------

\begin{itemize}
%
\item \bti{java\_object\_string/2}\\
\noindent\bt{java\_object\_string(ObjId, String)} is true iff
\bt{ObjId} is a term referencing a Java object and
\bt{PrologString} is the string representation of the object
(according to the semantics of the \texttt{toString} method
provided by the Java object).\\
\template{java\_object\_string(@obj\_id, ?string)}
%
\end{itemize}

%---------------------------------------------------------------------
\section{Functors}
%---------------------------------------------------------------------

No functors are provided by the \texttt{JavaLibrary} library.

%---------------------------------------------------------------------
\section{Operators}
%---------------------------------------------------------------------

\begin{table}[h]
    %
    \begin{center}{\small\tt
    \begin{tabular}{p{2cm}|p{1cm}|p{1cm}}\hline\hline
    Name & Assoc. & Prio. \\ \hline\hline
    <-   & xfx & 800\\
    returns     & xfx & 850 \\
    as   & xfx & 200\\
    .   & xfx & 600\\
    \hline\hline
    \end{tabular}
    }\end{center}
\end{table}

%\clearpage


%-----------------------------------------------------------------------
\section{Java Library Examples}
%-----------------------------------------------------------------------

The following examples are designed to show \texttt{JavaLibrary}'s
ease of use and flexibility.

%-------------------------------
\subsection{RMI Connection to a Remote Object}
%-------------------------------

Here we connect via RMI to a remote Java object.
%
In order to allow the reader to try this example with no need of
other objects, we connect to the remote Java object identified by
the name \verb|'prolog'|, which is an RMI server bundled with
the \tuprolog{} package, and can be spawned by typing:

{\small%
\texttt{java -Djava.security.all=policy.all  alice.tuprologx.runtime.rmi.Daemon}
}

\noindent Then, we invoke the object method whose signature is

{\small%
\texttt{SolveInfo solve(String goal);}
}
%
{\small%
\begin{verbatim}
    ?-  java_object('java.rmi.RMISecurityManager', [], Manager),
        class('java.lang.System') <- setSecurityManager(Manager),
        class('java.rmi.Naming') <- lookup('prolog') returns Engine,
        Engine <- solve('append([1],[2],X).') returns SolInfo,
        SolInfo <- success returns Ok,
        SolInfo <- getSubstitution returns Sub,
        Sub <- toString returns SubStr, write(SubStr), nl,
        SolInfo <- getSolution returns Sol,
        Sol <- toString returns SolStr, write(SolStr), nl.
\end{verbatim}
}
%
\noindent The Java version of the same code would be:
%
{\small%
\begin{verbatim}
        System.setSecurityManager(new RMISecurityManager());
        PrologRMI core = (PrologRMI) Naming.lookup("prolog");
        SolveInfo info = core.solve("append([1],[2],X).");
        boolean ok = info.success();
        String sub = info.getSubstiturion();
        System.out.println(sub);
        String sol = info.getSolution();
        System.out.println(sol);
\end{verbatim}
}


%-------------------------------
\subsection{Java Swing GUI from \tuprolog}
%-------------------------------

What about creating Java GUI components from the \tuprolog{}
environment?
%
Here is a little example, where a standard Java Swing open file
dialog windows is popped up:
%
{\small%
\begin{verbatim}
    open_file_dialog(FileName):-
        java_object('javax.swing.JFileChooser', [], Dialog ),
        Dialog <- showOpenDialog(_) returns Result,
        write(Result),
        Dialog <- getSelectedFile returns File,
        File <- getName returns FileName,
        class('java.lang.System') . out <- get(Out),
        Out <- println('you want to open file '),
        Out <- println(FileName).
\end{verbatim}
}

%-------------------------------
\subsection{Database access via JDBC from \tuprolog}
%-------------------------------

This example shows how to access a database via the Java standard
JDBC interface from \tuprolog{}.
%
The program computes the minimum path between two cities, fetching
the required data from the database called `distances'.
%
The entry point of the Prolog program is the \texttt{find\_path}
predicate.
%
{\small%
\begin{verbatim}
    find_path(From, To) :-
        init_dbase('jdbc:odbc:distances', Connection, '', ''),
        exec_query(Connection,
          'SELECT city_from, city_to, distance FROM distances.txt',
          ResultSet),
        assert_result(ResultSet),
        findall(pa(Length,L), paths(From,To,L,Length), PathList),
        current_prolog_flag(max_integer, Max),
        min_path(PathList, pa(Max,_), pa(MinLength,MinList)),
        outputResult(From, To, MinList, MinLength).

    paths(A, B, List, Length) :-
        path(A, B, List, Length, []).

    path(A, A, [], 0, _).
    path(A, B, [City|Cities], Length, VisitedCities) :-
        distance(A, City, Length1),
        not(member(City, VisitedCities)),
        path(City, B, Cities, Length2, [City|VisitedCities]),
        Length is Length1 + Length2.

    min_path([], X, X) :- !.
    min_path([pa(Length, List) | L],  pa(MinLen,MinList), Res) :-
        Length < MinLen, !,
        min_path(L, pa(Length,List), Res).
    min_path([_|MorePaths], CurrentMinPath, Res) :-
        min_path(MorePaths, CurrentMinPath, Res).

    writeList([]) :- !.
    writeList([X|L]) :- write(','), write(X), !, writeList(L).

    outputResult(From, To, [], _) :- !,
        write('no path found from '), write(From),
        write(' to '), write(To), nl.
    outputResult(From, To, MinList, MinLength) :-
        write('min path from '), write(From),
        write(' to '), write(To), write(': '),
        write(From), writeList(MinList),
        write('  - length: '), write(MinLength).

    % Access to Database

    init_dbase(DBase, Username, Password, Connection) :-
        class('java.lang.Class') <- forName('sun.jdbc.odbc.JdbcOdbcDriver'),
        class('java.sql.DriverManager') <- getConnection(DBase, Username, Password)
            returns Connection,
        write('[ Database '), write(DBase), write(' connected ]'), nl.

    exec_query(Connection, Query, ResultSet):-
        Connection <- createStatement returns Statement,
        Statement <- executeQuery(Query) returns ResultSet,
        write('[ query '), write(Query), write(' executed ]'), nl.

    assert_result(ResultSet) :-
        ResultSet <- next returns Valid, Valid == true, !,
        ResultSet <- getString('city_from') returns From,
        ResultSet <- getString('city_to') returns To,
        ResultSet <- getInt('distance') returns Dist,
        assert(distance(From, To, Dist)),
        assert_result(ResultSet).
    assert_result(_).
\end{verbatim}
}

%-------------------------------
\subsection{Dynamic compilation}
%-------------------------------

As already said, the \texttt{java\_class} predicate performs
\textit{dynamic compilation}, creating an instance of a Java
\texttt{Class} class that represents the public class declared in
the source text provided as argument.
%
The created \texttt{Class} instance, referenced by a Prolog term,
can be used to create instances via the \texttt{newInstance}
method, to retrieve specific constructors via the
\texttt{getConstructor} method, to analyze class methods and
fields, and for other above-mentioned meta-services: a sketch is
reported in \xf{dynamic-compilation}.
%
The \texttt{java\_class} arguments in the example specify, besides
the source text and the binding variable, the full class name
(\texttt{Counter}), which is necessary to locate the class in the
package hierarchy, and possibly a list of class paths required
for a successful compilation (if any).

\begin{figure}
\caption{Predicate \texttt{java\_class} performing dynamic compilation of Java code in \tuprolog{}.
\labelfig{dynamic-compilation}}
\begin{verbatim}
    ?- Source = 'public class Counter { ... }',
       java_class(Source, 'Counter', [], counterClass),
       counterClass <- newInstance returns myCounter,
       myCounter <- setValue(5),
       myCounter <- getValue returns X,
       write(X).
\end{verbatim}
\end{figure}

\xf{jcexamples} shows a more complex example, where a Java source
is retrieved via FTP and then exploited first to create a new
(previously unknown) class, and then a new instance of that class.
(The FTP service is provided by a shareware Java library.)
%
\begin{figure}
\caption{A new Java class is compiled and used after being retrieved via FTP.
\labelfig{jcexamples}}
{\scriptsize
\begin{verbatim}
% A user whose name is 'myName' and whose password is 'myPwd' gets the content of the file
% 'Counter.java' from the server whose IP address is 'srvAddr', creates the corresponding
% Java class and exploits it to instantiate and deploy an object

test :-
    get_remote_file('alice/tuprolog/test', 'Counter.java', srvAddr, myName, myPwd, Content),
    % creating the class
    java_class(Content, 'Counter', [], CounterClass),
    % instantiating (and using) an object of such a class
    CounterClass <- newInstance returns MyCounter,
    MyCounter <- setValue(303),
    MyCounter <- inc,
    MyCounter <- inc,
    MyCounter <- getValue returns Value,
    write(Value), nl.

% +DirName: Directory on the server where the file is located
% +FileName: Name of the file to be retrieved
% +FTPHost: IP address of the FTP server
% +FTPUser: User name of the FTP client
% +FTPPwd: Password of the FTP client
% -Content: Content of the retrieved file

get_remote_file(DirName, FileName, FTPHost, FTPUser, FTPPwd, Content) :-
    java_object('com.enterprisedt.net.ftp.FTPClient', [FTPHost], Client),
    % get file
    Client <- login(FTPUser, FTPPwd),
    Client <- chdir(DirName),
    Client <- get(FileName) returns Content,
    Client <- quit.
\end{verbatim}
}
\end{figure}
%
Though a lot remains to explore, \texttt{java\_class} features
seem quite interesting: in perspective one might think, for
instance, of a Prolog intelligent agent that dynamically acquires
information on a Java resource, and then autonomously builds up,
at run-time, the proper Java machinery enabling efficient
interaction with the resource.

%=======================================================================
\chapter{How to Develop New Libraries}
\label{ch:howto-develop-libraries}
%=======================================================================

Libraries are \tuprolog{}'s way to achieve the desired characteristics
of minimality, dynamic configurability, and straightforward
Prolog-to-Java integration.
%
Libraries are reflection-based, and can be written both in Prolog
and Java: other languages may be used indirectly, via JNI (Java
Native Interface).
%
At the \tuprolog{} side, exploiting a library written in Java
requires no pre-declaration of the new built-ins, nor any other
special mechanism: all is needed is the presence of the
corresponding \texttt{.class} library file in the proper location
in the file system.

\section{Implementation details}

Syntactically, a library developed in Java must extend the base
abstract class \texttt{alice.tuprolog.Library}, provided within
the \tuprolog{} package, and define new \textit{predicates} and/or
\textit{evaluable functors} and/or \textit{directives} in the form
of methods, following a simple signature convention.
%
In particular, new predicates must adhere to the signature:
%
\begin{center}
\small\tt
    public boolean <\textit{pred name}>\_<\textit{N}>(\textit{T1} arg1,
\textit{T2} arg2, ...,\textit{Tn} argN)
\end{center}
%
while evaluable functors must follow the form:
%
\begin{center}
    \small\tt
    public Term <\textit{eval funct name}>\_<\textit{N}>(\textit{T1} arg1,
\textit{T2} arg2, ...,\textit{Tn} argN)
\end{center}
%
and directives must be provided with the signature:
%
\begin{center}
    \small\tt
    public void <\textit{dir name}>\_<\textit{N}>(\textit{T1} arg1,
\textit{T2} arg2, ..., \textit{Tn} argN)
\end{center}
%
where \textit{T1}, \textit{T2}, ... \textit{Tn} are \texttt{Term} or derived
classes, such as \texttt{Struct}, \texttt{Var}, \texttt{Long}, etc., defined in
the \tuprolog{} package, constituting  the Java counterparts of
the corresponding Prolog data types.
%
The parameters represent the arguments actually passed to the built-in
predicate, functor, or directive.

%
A library defines also a new piece of theory, which is collected
by the Prolog engine through a call to the library method \texttt{String
getTheory()}.
%
By default, this method returns an empty theory: libraries which need to
add a Prolog theory must override it accordingly.
%
Note that only the external representation of a library's theory is
constrained to be in \texttt{String} form; the internal implementation
can be freely chosen by the library designer. However, using a Java
\texttt{String} for wrapping a library's Prolog code guarantees
self-containment while loading libraries through remote mechanisms such
as RMI.

\begin{table}[h]
    %
    \caption{Predicate and functor definitions in Java and their use
    in a \tuprolog{} program.\labeltab{java-preds}}
    %
    \begin{center}{\small\tt
    \begin{tabular}{p{10cm}|p{3.25cm}}
     \hline
     & \\
    \textit{// sample library} & \textit{\% tuProlog test program}\\
     import alice.tuprolog.*; & \\
     & \\
    public class TestLibrary extends Library \{                            &
 test :-\\
    ~~\textit{// builtin functor sum(A,B)}                          & ~~N is sum(5,6),\\
    ~~public Term sum\_2(Number arg0, Number arg1)\{    & ~~println(N).\\
    ~~~~float   arg0 = arg0.floatValue();        & \\
    ~~~~float   arg1 = arg1.floatValue();        & \\
    ~~~~return new Float(arg0+arg1);                                & \\
    ~~\}                                                            & \\
    ~~\textit{// builtin predicate println(Message)}                & \\
    ~~public boolean println\_1(Term arg)\{                      & \\
    ~~~~System.out.println(arg);                                   & \\
    ~~~~return true;                                                & \\
    ~~\}                                                            & \\
    \}                                                              & \\
     & \\
     \hline
    \end{tabular}
    }\end{center}
\end{table}

\xt{java-preds} shows a couple of examples about how a predicate
(such as \texttt{println/1}) and an evaluable functor (such as
\texttt{sum/2}) can be defined in Java and exploited from
\tuprolog{}.
%
The Java method \texttt{sum\_2}, which implements the evaluable
functor \texttt{sum/2}, is passed two \texttt{Number} terms (5 and 6)
which are then used (via \texttt{getTerm}) to retrieve the two
(float) arguments to be summed.
%
In the same way, method \texttt{println\_1}, which implements the
predicate \texttt{println/1}, receives \texttt{N} as \texttt{arg},
and retrieves its actual value via \texttt{getTerm}: since this is
a predicate, a boolean value (\texttt{true} in this case) is returned.
%

The developer of a library may face two corner case as far as method
naming is concerned: the first happens when the name of the
predicate, functor or directive she is defining contains a symbol
which cannot legally appear in a Java method's name; the second
occurs when he has to define a predicate and a directive with the
same Prolog signature, which Java would not be able to tell apart
because it cannot distinguish signatures of methods differing for
their return type only.
%
To overcome this kind of issues, a {\em synonym map} can be
constructed under the form of an array of \texttt{String} arrays,
and returned by the appropriate \texttt{getSynonymMap} method,
defined as abstract by the \texttt{Library} class. In both the cases
described above, another name must be chosen for the Prolog
executable element the library's developer want to define: then, by
means of the synonym map, that fake name can be associated with the
real name and the type of the element, be it a predicate, a functor
or a directive.
%
For example, if a definition for an evaluable functor representing
addition is needed, but the symbol \texttt{+} cannot appear in a
Java method's name, a method called \texttt{add} can be defined and
associated to its original Prolog name and its function by inserting
the array \texttt{\{"+", "add", "functor"\}} in the synonym map.

\section{Library Name}

% Lib name
%
By default, the name of the library coincides with the full class name of the
class implementing it.
%
However, it is possible to define explicitly the name of a library by
overriding the \texttt{getName} method, and returning as a string
the real name.
%
For example:
%
\begin{verbatim}
package acme;
import alice.tuprolog.*;
public class MyLib_ver00 extents Library {
    public String getName(){
        return "MyLibrary";
    }
    ...
}
\end{verbatim}

This class defines a library called \texttt{MyLibrary}.
%
It can be loaded into a Prolog engine by using the
\texttt{loadLibrary} method on the Java side, or a
\texttt{load\_library} built-in predicate on the Prolog side,
specifying the full class name (\texttt{acme.MyLib\_{ve00}}).
%
It can be unloaded then dynamically using the \texttt{unloadLibrary}
method (or the corresponding \texttt{unload\_library} built-in),
specifying instead the \textit{library name} (\texttt{MyLibrary}).

%\include{developerGuideAppendix}
\bibliography{bibliography/bibliography}
\bibliographystyle{plain}

\end{document}
